\scanpage{153}\mspage{27}
\unsure{en|un|$\sim$}
\[
C_1=\mathrm{card}\,\underbrace{\mathrm{Fl}\,M^{\pi_1}}=\mathrm{card}\,\mathrm{Fl}^{\pi_1}M
\]
(où pour tt $\pi\ge\pi_0$, $M^\pi$ désigne la sous-catégorie \inserted{\unsure{str.|str}} pleine de $M$ définie par $\mathrm{Filt}^\pi\deleted{\unsure{)|\}}}$, et $\mathrm{Fl}^\pi M=\mathrm{Fl}\,M^\pi$).
Soit $C_0$ une petite partie de $\mathrm{Fl}\,M$, formée de flèches \unsure{coquarrables}, et posons
\[
\Phi=\widetilde{C}_0=\Psi_{\deleted{\unsure{*|\alpha}}}^{*},\quad\text{\unsure{où|i.e.}}\quad\Psi=C_{0\unsure{\alpha|*}}=\Phi_*.
\]
Soit $r$ un ordinal $\ge\mathrm{Sup}(\pi_1,C_1)$, et tel que l'on ait $C_0\subset\mathrm{Fl}^rM$. On suppose que $C_0$ est irrédundant (i.e. deux objets de $C_0$ isomorphes sont égaux) plus généralement, que $C_0$ satisfait la \inserted{très anodine} condition de non pléthoricité (3.11.) relativement à $r$.
On considère la suite transfinie des foncteurs $p_\alpha=p_\alpha^{C_0}$ sur $\underline{\mathrm{Fl}}\,M$ (cf. \unsure{N.|M.}1), \deleted{Alors} \deleted{et \unsure{on} rappelle} \deleted{que \unsure{ce} qui} donnant \unsure{lieu|une} \unsure{aux|des} factorisations fonctorielles d'une $f\in\underline{\mathrm{Fl}}\,M$, $f:X\mathrel{\text{---}}Y$, en
\diagram{
\begin{tikzcd}[ampersand replacement=\&]
\unsure{x|X}\arrow[r,"{i_\alpha(f)}"]\&X_\alpha(f)\arrow[r,dash,"{p_\alpha(f)}","{\unsure{\phantom{\longrightarrow}}}"']\&Y.
\end{tikzcd}
\quad\unsure{(3.16,|(3.16)|(3.16.)}
}
On a alors ce qui suit :

1°) On a (\unsure{bien entendu|bien sûr}) $i_\alpha(f)\in\Phi$ pour tout $\alpha$, et de plus, \uline{si $\alpha$ est grand devant $r^{\pi_0}$}, on a aussi $p_\alpha(f)\in\Psi$. Si $r'$ est un cardinal $>r^{\pi_0}$ (p. ex. le \deleted{\unsure{son|le|un}} cardinal successeur de $r^{\pi_0}$), on peut prendre pour $\alpha$ le premier ordinal dont le cardinal soit égal à $r'$.

2°) \deleted{Supposons $\alpha$ choisi grand devant $\unsure{r^{\pi_0}|\pi_0}$, Alors} Pour tout cardinal $s$ de la forme
\[
s=c^t\quad\text{avec}\quad t\ge r^{\pi_0},\quad c\ge2
\]
(on a $s=2^t$ si on prend $c\le t\deleted{\unsure{)|\}}}$, cf. lemme \unsure{3.2|3.1}), $\underline{\mathrm{Fl}}^sM$ est stable sous les $p_\alpha$ tels que $\mathrm{card}\,\alpha\le s$. Si on prend p.ex. pour $\alpha$ le premier ordinal dont le cardinal soit $>r^{\pi_0}$, alors $\underline{\mathrm{Fl}}^sM$ est stable par $p_\alpha^{\deleted{\unsure{C_0|*|r'}}}$, qui par 1°) réalise donc une \unsure{factorisation fonctorielle|\deleted{factorisation fonctorielle}} (3.16) avec $i_\alpha(f)\in\Phi$, $p_\alpha(f)\in\Psi$.
% archive sheet number: 152
